\section{INTRODUCTION}
    \subsection{Background of the project}
            Armed Forces and Police Mutual Benefit Association Incorporated or AFPMBAI manages a large volume of documents related to Board of Trustees like meetings, board resolutions and administrative records. These documents are currently stored in fragmented formats and locations like google drive for quick collaboration, local computer storage for draft versions and an on-premise server storage for historical records which makes the document retrieval time consuming and inefficient. The absence of a centralized and intelligent document archiving system often resulted in redundant searches, late decision-making and limited accessibility to institutional knowledge.

Though AFPMBAI operates an on-premise server for document records, there are certain files remain in cloud-based storage applications like Google Drive for convenience and accessibility of trustees and the Board Relations Office (BRO) staff. The multiple storage of documents approach might likely be practical but it has a risks of data redundancy, data inconsistency and not complying with the organization’s strict data privacy policy. AFPMBAI requires that all official systems, particularly those containing Board records, be housed on-site and managed internally. This is to ensure complete control over data security, given the institution's cautious approach and the very sensitive information it manages.

Currently, when management needs documents from the Board Relations Office (BRO), staff has to search for manually using simple keyword searches on their documents. They often make a table of contents for each document to help find it later, but this isn't practical anymore because there are too many documents to look through. Another issue is the difficulty of connecting older information. Employees spend several hours each week looking through folders and different file versions to find important documents. This slows down report writing and decision-making. 

In contrast, many organizations in the financial, government, and defense sectors have successfully implemented intelligent document archiving systems powered by artificial intelligence (AI). For example, the Commonwealth Bank of Australia integrated H2O.ai’s Document AI platform is to streamline the processing of both structured and unstructured data across billions of transactions (Pratusevich, 2024). Within four months, the bank was able to process millions of documents per day accelerating customer onboarding and strict compliance with regulatory and risk policies. The system automatically extracted information like names, birth dates, and addresses from identification documents which enables faster verification while maintaining accuracy. The result of the bank achieved up to 10 times faster invoice processing and 50\%- 85\% automation accuracy across different document types. This case shows that using AI in document management can make operations more efficient, help meet regulations, and speed up data processing without risking the privacy and control of sensitive information.

As outlined, we need to create an on-premise smart document archiving system. This system included a locally integrated large language model (LLM) to allow secure, context-aware searching and chatbot-supported interaction with documents. The goal is to simplify document archiving and improve how users experience it, letting them ask questions in natural language without risking data privacy.


    \subsection{Objectives}
        The study aims to create an AI-supported document archiving system for the AFPMBAI Board of Trustees. This system securely handle and manage board documents using an on-premise platform that includes an AI-driven chatbot.

                \begin{itemize}
                \item Implement a central location to safely keep and organize all Board of Trustees documents, such as meeting minutes, decisions, rules, presentations, and memos.
                \begin{itemize}
                    \item  Data is stored encrypted with AES-256 to keep the data confidential.
                    \item The system can only be used within AFPMBAI offices or remotely via a VPN set up that managed by the AFPMBAI Infrastructure Team.
                    \item Control who can access the system using role-based access control (RBAC) and keep audit trails for extra security.
                \end{itemize}
                \item Integrate an on-premise Large Language Model (LLM) that is capable of processing textual data and enabling natural language queries without depending on cloud-based AI services.
                \item Develop an intelligent search and chatbot interface that allows users to retrieve information contextually through context-aware interaction.
                 \item Evaluate the performance of the AI-assisted search mechanism based on its ability to retrieve relevant board documents compared to the existing manual retrieval process.

\end{itemize}

    \subsection{Scope and Limitations}
        \subsubsection{Scope}
           This study focused on designing and building an AI-Assisted document archiving system for the Board of Trustees (BOT) of the Armed Forces and Police Mutual Benefit Association Incorporated (AFPMBAI). The system's goal is to offer a secure and smart on-premise application for managing and finding board-related documents, such as meeting minutes, resolutions, memoranda, and policy documents.The system uses a central repository on AFPMBAI's own servers to consolidate all board documents into one secure location. Additionally, it utilized AES-256 encryption to ensure stored data stays private and only authorized users can access it. The system also incorporates the Docling library for parsing documents and OCR which allows it to extract, organize, and prepare text from various document types for efficient processing and retrieval. This processed data are utilized by a locally hosted Large Language Model, specifically Gemma 3:1B, which is selected for its efficiency and ability to operate within constrained on-premise server resources. The model allowed users to ask questions naturally and get answers that consider the conversation's context. The project has a web interface available only within the organization's network, with different access levels for administrators, BRO staff, and board of trustees. It measures the system's success by how accurate the chatbot is, how fast it finds information, and how much it cuts down the time people spend searching for documents compared to our current manual process. The project used an Agile development approach, which means the building and delivering modules of the system in stages, working closely with users involved. I organize development into short cycles (sprints), and each cycle includes refining needs, designing, building, testing, and getting user feedback. This way of working improves the system gradually, checks that features work early on, and adjusts to changing user needs. The study will result in a complete, working system for archiving board of trustees documents, supported by AI along with all necessary technical documentation.
           
        \subsubsection{Limitations}  
            The study is subject to several limitations as this project will only focus on document management and intelligent retrieval functions which excludes the broader record management capabilities like physical document tracking, workflow approvals, and records lifecycle management. The AI chatbot and language model will operate entirely on-premise in compliance with AFPMBAI’s data privacy policies which have no integration with cloud-based AI platforms to be implemented in which it may limit the model’s sophistication. The accuracy and contextual understanding of the AI model will depend on the quality and quantity of existing board documents available for training, restricted to pre-approved internal materials. Additionally, the system’s performance may vary depending on the available server resources such as CPU, RAM, and storage capacity within AFPMBAI’s local infrastructure. The development will focus on delivering a functional and iterative system through Agile methodology rather than a fully scaled enterprise deployment. Additionally, some advanced features, optimizations, and integrations may not be fully realized within the study period. Despite these limitations, the study aims to demonstrate the feasibility, security, and operational efficiency of implementing an AI-powered, on-premise document archiving system within an organization operating in a conservative and security-sensitive environment. \\
            